\emph{Preliminaries.} By Theorem \claimref{std-algebra.witt-basis-car} the $v_{j}, w_{j}$ satisfy the canonical anticommutation relations, $v_{j}^{2} = w_{j}^{2} = 0$, $\{v_{j}, w_{k}\} = \delta_{jk}$, and distinct planes' generators anticommute; $n_{j} = v_{j}w_{j} = \tfrac12(1 - e_{j}f_{j})$ is an idempotent, $n_{j}w_{j} = 0$, $w_{j}n_{j} = w_{j}$, $n_{j}v_{j} = v_{j}$, $v_{j}n_{j} = 0$. $J^{2} = g_{1}g_{2}g_{1}g_{2} = -g_{1}^{2}g_{2}^{2} = -1$, and $J$ commutes with every $e_{j}$, $f_{j}$ (each anticommutes with $g_{1}$ and with $g_{2}$). The adjoint $*$ is an antiautomorphism with $e_{j}^{*} = e_{j}$, $f_{j}^{*} = -f_{j}$, $g_{a}^{*} = -g_{a}$, so $v_{j}^{*} = w_{j}$ and $J^{*} = g_{2}^{*}g_{1}^{*} = g_{2}g_{1} = -J$. For a basis blade $u$ (a product of distinct generators) $u^{*}u = 1$: each generator $g$ has $g^{*}g = 1$ ($g^{*} = g$ when $Q(g) = 1$, $g^{*} = -g$ when $Q(g) = -1$), and the reversion in $*$ pairs each generator with itself from the inside out. Two facts from the preliminaries of Theorem \claimref{synthesis.fluctuation-forms-of-the-tower-and-density-wave-bosons-of-the-null-vacuum} are used: elements of even degree on disjoint sets of planes commute, and the scalar part factorises over disjoint sets of planes, so a product state $\omega_{N}$ is multiplicative on elements of disjoint blocks; extended $\mathrm{span}\{1, J\}$-linearly, both remain true, since $J$ commutes with the transport algebra and every product expands into terms $J^{a}$ times products of transport elements.

\emph{(i)} Generators of disjoint planes anticommute, and $J$ commutes with all of them, so moving a product of $a$ such generators past a product of $b$ others gives the sign $(-1)^{ab}$; by linearity $xy = yx$ unless $x$ and $y$ are both odd, when $xy = -yx$ (graded locality, clause (ii) of \claimref{synthesis.ladder-bosons-inherit-spin-statistics-from-the-single-modes}). Then $[x_{k}, y_{l}]_{\epsilon} = 0$ for $k \ne l$ with $\epsilon$ as stated, which gives the displayed identity; for odd $x$, $y$ the commutator of $x_{k}$ and $y_{l}$ with $k \ne l$ is $2x_{k}y_{l}$, in general nonzero (the check exhibits it for $x = w$, $y = v$). In an even product state, $\omega(x) = 0$ for odd $x$, so for odd elements $F_{N}(x) = N^{-1/2}\sum_{k}x_{k}$; in all cases $[F_{N}(x), F_{N}(y)]_{\epsilon} = N^{-1}\sum_{k}C_{k}$ ($C$ even), and the mean and variance follow as in Theorem \claimref{synthesis.fluctuation-forms-of-the-tower-and-density-wave-bosons-of-the-null-vacuum}(vii): the $C_{k} - \omega(C)$ are independent and centred, so the cross terms of the variance vanish and each block contributes $\omega(C^{*}C) - \omega(C^{*})\omega(C)$. The trace: $\br{xy}_{0} = \br{yx}_{0}$ for all elements of a Clifford algebra, and $\br{Jxy}_{0} = \br{yJx}_{0} = \br{Jyx}_{0}$ for $x, y \in \mathcal{T}_{J}$ because $J$ commutes with $y$; so $\mathrm{tr}([x, y]) = 0$, while $\mathrm{tr}(\alpha + \beta J) = \alpha + \beta J$ for real $\alpha, \beta$ ($\br{J}_{0} = 0$, $\br{J^{2}}_{0} = -1$). A commutator of elements of $\mathcal{T}_{J}$ equal to a nonzero element of $\mathrm{span}\{1, J\}$ would have nonzero trace. The restriction to $\mathcal{T}_{J}$ is needed: the cyclicity uses that $J$ commutes with $y$, and in the whole level $[g_{1}, g_{2}] = 2J$ (the check computes it). This is the trace argument of Proposition \claimref{field-theory.no-canonical-commutation-relations-tower} over the field $\mathrm{span}\{1, J\}$. The check verifies graded locality on two and three blocks for odd elements ($w$, the cubic $w_{1}n_{2}$, $v_{1} + v_{1}v_{2}w_{2}$) and even ones (hopping, pairing, $\phi + J\chi$), the bracket identities, and the mean $p_{2}$ and variance $p_{2}(1 - p_{2})/N$ of the odd bracket $\{F_{N}(x), F_{N}(x)^{*}\}$, $x = w_{1}n_{2}$, $C = n_{2}$, in the product state $(\tfrac34, \tfrac15)$.

\emph{(ii)} Positivity of the extension: for $z = X + JY$ with $X, Y \in \mathcal{A}_{e}$, $z^{*}z = X^{*}X + Y^{*}Y + J(X^{*}Y - Y^{*}X)$, and $\omega(Y^{*}X) = \omega((X^{*}Y)^{*}) = \omega(X^{*}Y)$ for the real state, so $\omega(z^{*}z) = \omega(X^{*}X) + \omega(Y^{*}Y) \ge 0$. The form $(z, z') \mapsto \omega(z^{*}z')$ is therefore a positive sesquilinear form over the field $\mathrm{span}\{1, J\}$, with the Cauchy-Schwarz inequality; for a unitary $u$, $|\omega(u)|^{2} \le \omega(1)\omega(u^{*}u) = 1$. Every blade is unitary, so $|\omega(x)| \le \|x\|_{1}$, and $\|xy\|_{1} \le \|x\|_{1}\|y\|_{1}$ because a product of blades is $\pm$ a blade ($\|Jx\|_{1} = \|x\|_{1}$). For self-adjoint $y$, $(Jy)^{*} = -Jy$, so $e^{Jy}$ is unitary. Since $\omega(x^{*})$ is the $J$-conjugate of $\omega(x)$, $\omega$ is real on self-adjoint elements and in $J\R$ on anti-self-adjoint ones: $s$ is real and $\sigma_{\omega}(x, y) = \omega([x, y]) \in J\R$.

(a) $F_{N}(x_{j}) = N^{-1/2}\sum_{k}y_{j,k}$ and the $y_{j,k}$ of distinct blocks commute (even elements on disjoint planes, $J$ commuting), so $W_{N}(x_{j}) = \prod_{k}e^{Jy_{j,k}/\sqrt N}$ and the product over $j$ can be regrouped block by block, keeping the order of $j$ within each block: $\prod_{j}W_{N}(x_{j}) = \prod_{k}Z_{k}$, $Z_{k}$ the copy on block $k$ of $Z = e^{Jy_{1}/\sqrt N}\cdots e^{Jy_{r}/\sqrt N}$. Multiplicativity of $\omega_{N}$ on disjoint blocks gives $[\omega(Z)]^{N}$.

(b) Let $X_{j} = \sum_{k}y_{j,k}$. For one block added to $N - 1$ others, with $A_{j}$ the part on the first $N - 1$ blocks and $B_{j}$ the copy on the new block, $A$'s commute with $B$'s, so for an ordered list $T$ of indices $\omega_{N}(\prod_{t \in T}(A_{t} + B_{t})) = \sum_{U \subseteq T}\omega_{N}(\prod_{U}A)\,\omega_{N}(\prod_{T \setminus U}B)$, each product in the original order. In the ring of multilinear polynomials in commuting variables $t_{1}, \dots, t_{k}$ with $t_{j}^{2} = 0$, put $M(t) = \sum_{T}m(T)t^{T}$ for the ordered moments and $K(t) = \sum_{T \ne \emptyset}\kappa(T)t^{T}$; the moment-cumulant relation over set partitions is $M = \exp K$. The subset convolution is $M_{A + B} = M_{A}M_{B}$, so $K_{A + B} = K_{A} + K_{B}$, and by induction the cumulants of the sums $X_{j}$ are $N$ times those of one block; the normalisation $N^{-1/2}$ per entry gives $N^{1 - k/2}$, and first cumulants vanish as $\omega(y_{j}) = 0$.

(c) Let $z_{j} = Jy_{j}/\sqrt N$, $\|z_{j}\|_{1} = \|y_{j}\|_{1}/\sqrt N$. Expanding the exponentials, $Z = \sum_{d \ge 0}Z_{d}$ with $Z_{d}$ the part of total degree $d$ in the $z_{j}$ and $\|Z_{d}\|_{1} \le (\nu/\sqrt N)^{d}/d!$. $Z_{0} = 1$; $\omega(Z_{1}) = 0$; $Z_{2} = \sum_{j < l}z_{j}z_{l} + \tfrac12\sum_{j}z_{j}^{2} = -N^{-1}\bigl(\tfrac12(\sum_{j}y_{j})^{2} + \tfrac12\sum_{j < l}[y_{j}, y_{l}]\bigr)$, so $\omega(Z_{2}) = -c/N$. Hence $a := \omega(Z) = 1 - c/N + \varepsilon$ with $|\varepsilon| \le \sum_{d \ge 3}(\nu/\sqrt N)^{d}/d! \le \tfrac16(\nu/\sqrt N)^{3}e^{\nu/\sqrt N}$. $|a| \le 1$ ($Z$ is unitary), and $b := e^{-c/N}$ has $|b| = e^{-\operatorname{Re}c/N} \le 1$ because $\operatorname{Re}c = \tfrac12\omega((\sum_{j}y_{j})^{2}) \ge 0$, the commutator terms lying in $J\R$. Then $|a^{N} - b^{N}| \le N|a - b|$, and $|a - b| \le |1 - c/N - e^{-c/N}| + |\varepsilon| \le \tfrac12(|c|/N)^{2}e^{|c|/N} + |\varepsilon|$ with $|c| \le \tfrac12 \nu^{2}$ (the element whose value is $c$ has $\|\cdot\|_{1} \le \tfrac12(\sum_{j}\|y_{j}\|_{1})^{2}$). Multiplying by $N$ gives the bound, since $a^{N}$ is the left side of (a) and $b^{N} = e^{-c}$.

(d) In the Weyl algebra, $W(x_{1})\cdots W(x_{r}) = e^{-\frac12\sum_{j < l}\sigma_{\omega}(x_{j}, x_{l})}W(\sum_{j}x_{j})$ by induction on the Weyl relation, and $\omega_{\infty}$ of it is $\exp(-\tfrac12\sum_{j < l}\sigma_{\omega}(x_{j}, x_{l}) - \tfrac12 s(\sum x_{j}, \sum x_{j}))$, which is $e^{-c}$ because $s(\sum x_{j}, \sum x_{j}) = \omega((\sum_{j}y_{j})^{2})$ and $\sigma_{\omega}(x_{j}, x_{l}) = \omega([y_{j}, y_{l}])$. The finite-$N$ operators satisfy $W_{N}(x)^{*} = W_{N}(-x)$, so for $\gamma_{j} \in \mathrm{span}\{1, J\}$ the limit of $\omega_{N}(Z_{N}^{*}Z_{N}) \ge 0$, $Z_{N} = \sum_{j}\gamma_{j}W_{N}(x_{j})$, is $\sum_{j, l}\bar\gamma_{j}\gamma_{l}\,\omega_{\infty}(W(x_{j})^{*}W(x_{l}))$, which is therefore nonnegative: the limit functional is positive. Cauchy-Schwarz gives $|\omega(y_{x}y_{y})|^{2} \le \omega(y_{x}^{2})\omega(y_{y}^{2})$, and $\omega(y_{x}y_{y}) = s(x, y) + \tfrac12\sigma_{\omega}(x, y)$ with the first term real and the second in $J\R$, so $s(x, y)^{2} + \tfrac14|\sigma_{\omega}(x, y)|^{2} \le s(x, x)s(y, y)$. For the hopping fields, $\omega(b) = 0$ in a gauge-invariant product state, $\phi^{2} = bb^{*} + b^{*}b + b^{2} + b^{*2}$ with $b^{2} = 0$, $bb^{*} = n_{1}(1 - n_{2})$, $b^{*}b = n_{2}(1 - n_{1})$, so $s(\phi, \phi) = p_{1}(1 - p_{2}) + p_{2}(1 - p_{1})$; $\pi^{2} = -(b^{*} - b)^{2} = bb^{*} + b^{*}b$ gives the same; and $[\phi, \pi] = 2J[b, b^{*}] = 2J(n_{1} - n_{2})$ (Theorem \claimref{synthesis.fluctuation-forms-of-the-tower-and-density-wave-bosons-of-the-null-vacuum}(iii) for $[b, b^{*}]$); for a real state and self-adjoint $x, y \in \mathcal{A}_{e}$, $[x, y]^{*} = -[x, y]$ gives $\omega([x, y]) = -\omega([x, y])$, so the form vanishes there, as stored, of value $2J(p_{1} - p_{2})$. The third cumulant of the Bernoulli element $n_{1}$ is $p_{1} - 3p_{1}^{2} + 2p_{1}^{3} = p_{1}(1 - p_{1})(1 - 2p_{1})$. The rate: for one self-adjoint $x$ the element $y$ commutes with itself, so $\log\omega(e^{Jty}) = \sum_{k \ge 2}\kappa_{k}(Jt)^{k}/k! = -\kappa_{2}t^{2}/2 - J\kappa_{3}t^{3}/6 + O(t^{4})$ with $\kappa_{k}$ the cumulants of $y$, and by (a) $\omega_{N}(W_{N}(x)) = \exp\bigl(N\log\omega(e^{Jy/\sqrt N})\bigr) = e^{-c}\bigl(1 - J\kappa_{3}/(6\sqrt N) + O(N^{-1})\bigr)$, $c = \kappa_{2}/2$; a nonzero $\kappa_{3}$ makes the error exactly of order $N^{-1/2}$, with leading coefficient $e^{-c}|\kappa_{3}|/6$. The check computes the ordered cumulants of orders $2$ to $4$ for six ordered lists of the fields $\phi$, $\pi$, $n_{1}$, $\chi = c + c^{*}$ and $\phi + 2n_{1}n_{2}$ on $N = 1, 2, 3$ blocks in the product state $(\tfrac34, \tfrac15)$ and confirms the factor $N$ exactly; it confirms positivity of the extension, unitarity of all $64$ blades of $\Cl(2,4)$ and the uncertainty relation in three product states.

\emph{(iii)} In the polarised state every $b_{k}$ annihilates the state's vector, so normally ordered moments vanish. Let $K = \sqrt N B$ and $\psi_{l} = (K^{*})^{l}\Omega_{\mathrm{pol}}$. By Theorem \claimref{synthesis.fluctuation-forms-of-the-tower-and-density-wave-bosons-of-the-null-vacuum}(ii), $[K, K^{*}] = N[B, B^{*}]$ acts as $N - 2j$ on $\psi_{j}$, and $K\psi_{0} = 0$. Telescoping $K\psi_{j + 1} = K^{*}K\psi_{j} + (N - 2j)\psi_{j}$ over $j \le l$ gives $K\psi_{l + 1} = \sum_{j \le l}(N - 2j)\,\psi_{l} = (l + 1)(N - l)\psi_{l}$ (the ladder coefficient, also in \claimref{synthesis.ladder-states-are-locally-product-with-non-gaussian-fluctuations}(i) and in the quasi-spin of Proposition \claimref{field-theory.how-close-tower-comes-ccr}), so $\omega_{N}(B^{k}B^{*k}) = N^{-k}\prod_{l < k}(l + 1)(N - l) = k!\prod_{l < k}(1 - l/N)$. For the Gaussian limit, $s(\phi, \phi) = s(\pi, \pi) = 1$ and $\sigma_{\omega}(\phi, \pi) = 2J$ make the uncertainty bound an equality, and the limit of the moments $\omega_{N}(B^{k}B^{*k}) \to k!$, $\omega_{N}(B^{*k}B^{k}) = 0$ are those of the vacuum of a mode with $[a, a^{*}] = 1$: on the hopping pair $(\phi, \pi)$ the limit is the Fock vacuum (the other block fields, such as $n_{j}$ and $\chi$, are deterministic in the polarised state). The check computes these moments on $N = 1, 2, 3$ polarised blocks for $k = 1, \dots, N + 1$, and the third and fourth cumulants $0$ and $-2N$ of the unnormalised sum $\sum_{k}(b_{k} + b_{k}^{*})$.
